% GENERATED FILE. Edit canonical lessons, then run npm run generate:books.
% canonical-source-hash: fnv1a64:c75c861eb99f197c
% canonical-lessons: RU-C510-porazheniye, RU-C510-nichya, RU-C510-sopernik, RU-C510-proigrysh, RU-C510-lidirovat

\chapter{Winning and Losing}
\label{ch:ru-winning-and-losing}
\hlchaptermodality{510}

\begin{chapteropening}
\textbf{By the end of this chapter:} \emph{I can talk about a defeat, a tie and the lead.}

Complete the last lesson of chapter 510: I can talk about a defeat, a tie and the lead.
\end{chapteropening}

\section[\texorpdfstring{\ru{поражение} (porazhéniye)}{поражение (porazhéniye)}]{\ru{поражение} (porazhéniye) — a defeat}

\label{lesson:RU-C510-porazheniye}



\begin{quote}
Before the new one: say the Russian for a goal, then the Russian for to score.
\end{quote}

\subsection*{You'll want to know: \ru{поражение}}
\textbf{\ru{поражение}} (\emph{porazhéniye}) — \textquotedblleft{}a defeat\textquotedblright{}.

\begin{quote}
\ru{После} \ru{поражения} \ru{игроки} \ru{долго} \ru{молчали}.
\end{quote}

\begin{quote}
\emph{Pósle porazhéniya igrokí dólgo molcháli.}
\end{quote}

\begin{quote}
After the defeat the players were silent for a long time.
\end{quote}

Neuter. \emph{\ru{Потерпеть} \ru{поражение}} — to suffer a defeat.

\begin{cousinweb}
\emph{\ru{поражение}} is built on \emph{\ru{поразить}} 'to strike, to defeat', from \emph{\ru{разить}} 'to strike'; EtymDB lists it, glossed 'defeat', without a source.
\end{cousinweb}

\begin{grammarlens}[title={What you've built}]
One more word for telling, explaining and describing at B1.
\end{grammarlens}

\subsection*{Your turn}
\begin{itemize}
\raggedright
  \item \emph{Say it:} \emph{porazhéniye}
  \item \emph{Say it:} \emph{porazhéniye}, once more
  \item \emph{Say it:} a sentence of your own with \emph{porazhéniye}
  \item \emph{Recall:} say the Russian for a goal, then the Russian for to score, then say \emph{porazhéniye} again
\end{itemize}

\begin{tcolorbox}[breakable,colback=teal!4,colframe=teal!35!black,title={Before you move on}]
What does \ru{поражение} mean? (\textquotedblleft{}A defeat\textquotedblright{}.) Say it once more.
\end{tcolorbox}
\section[\texorpdfstring{\ru{ничья} (nichyá)}{ничья (nichyá)}]{\ru{ничья} (nichyá) — a tie (in a game)}

\label{lesson:RU-C510-nichya}



\begin{quote}
Before the new one: say the Russian for to score, then the Russian for a defeat.
\end{quote}

\subsection*{You'll want to know: \ru{ничья}}
\textbf{\ru{ничья}} (\emph{nichyá}) — \textquotedblleft{}a tie (in a game)\textquotedblright{}.

\begin{quote}
\ru{Ничья} \ru{устроила} \ru{обе} \ru{команды}.
\end{quote}

\begin{quote}
\emph{Nichyá ustróila óbe komándy.}
\end{quote}

\begin{quote}
The tie suited both teams.
\end{quote}

Feminine, declined like an adjective. \emph{\ru{Сыграть} \ru{вничью}} — to end level.

\begin{cousinweb}
\emph{\ru{ничья}} comes from \emph{\ru{ничья} \ru{игра}} 'nobody's game' (EtymDB), built on \emph{\ru{ничей}} 'nobody's'.
\end{cousinweb}

\begin{grammarlens}[title={What you've built}]
One more word for telling, explaining and describing at B1.
\end{grammarlens}

\subsection*{Your turn}
\begin{itemize}
\raggedright
  \item \emph{Say it:} \emph{nichyá}
  \item \emph{Say it:} \emph{nichyá}, once more
  \item \emph{Say it:} a sentence of your own with \emph{nichyá}
  \item \emph{Recall:} say the Russian for to score, then the Russian for a defeat, then say \emph{nichyá} again
\end{itemize}

\begin{tcolorbox}[breakable,colback=teal!4,colframe=teal!35!black,title={Before you move on}]
What does \ru{ничья} mean? (\textquotedblleft{}A tie (in a game)\textquotedblright{}.) Say it once more.
\end{tcolorbox}
\section[\texorpdfstring{\ru{соперник} (sopérnik)}{соперник (sopérnik)}]{\ru{соперник} (sopérnik) — a rival, an opponent}

\label{lesson:RU-C510-sopernik}



\begin{quote}
Before the new one: say the Russian for a defeat, then the Russian for a tie.
\end{quote}

\subsection*{You'll want to know: \ru{соперник}}
\textbf{\ru{соперник}} (\emph{sopérnik}) — \textquotedblleft{}a rival, an opponent\textquotedblright{}.

\begin{quote}
\ru{Сегодня} \ru{у} \ru{нас} \ru{очень} \ru{сильный} \ru{соперник}.
\end{quote}

\begin{quote}
\emph{Segódnya u nas óchen sílnyy sopérnik.}
\end{quote}

\begin{quote}
Today we have a very strong opponent.
\end{quote}

Masculine; feminine \ru{соперница}.

\begin{cousinweb}
\emph{\ru{соперник}} is inherited from Old East Slavic \emph{s\ru{ъ}p\ru{ь}r\ru{ь}nik\ru{ъ}}, built with the prefix that EtymDB traces to Proto-Slavic \emph{sǫ-} 'together'.
\end{cousinweb}

\begin{grammarlens}[title={What you've built}]
One more word for telling, explaining and describing at B1.
\end{grammarlens}

\subsection*{Your turn}
\begin{itemize}
\raggedright
  \item \emph{Say it:} \emph{sopérnik}
  \item \emph{Say it:} \emph{sopérnik}, once more
  \item \emph{Say it:} a sentence of your own with \emph{sopérnik}
  \item \emph{Recall:} say the Russian for a defeat, then the Russian for a tie, then say \emph{sopérnik} again
\end{itemize}

\begin{tcolorbox}[breakable,colback=teal!4,colframe=teal!35!black,title={Before you move on}]
What does \ru{соперник} mean? (\textquotedblleft{}A rival, an opponent\textquotedblright{}.) Say it once more.
\end{tcolorbox}
\section[\texorpdfstring{\ru{проигрыш} (próigrysh)}{проигрыш (próigrysh)}]{\ru{проигрыш} (próigrysh) — a loss (in a game)}

\label{lesson:RU-C510-proigrysh}



\begin{quote}
Before the new one: say the Russian for a tie, then the Russian for a rival.
\end{quote}

\subsection*{You'll want to know: \ru{проигрыш}}
\textbf{\ru{проигрыш}} (\emph{próigrysh}) — \textquotedblleft{}a loss (in a game)\textquotedblright{}.

\begin{quote}
\ru{Проигрыш} \ru{не} \ru{испортил} \ru{нам} \ru{настроение}.
\end{quote}

\begin{quote}
\emph{Próigrysh ne ispórtil nam nastroyéniye.}
\end{quote}

\begin{quote}
Losing didn't spoil our mood.
\end{quote}

Masculine. Built on \emph{\ru{проиграть}} — to lose.

\begin{cousinweb}
\emph{\ru{проигрыш}} is built on \emph{\ru{проиграть}} 'to lose', from \emph{\ru{игра}} 'a game'; EtymDB has no entry.
\end{cousinweb}

\begin{grammarlens}[title={What you've built}]
One more word for telling, explaining and describing at B1.
\end{grammarlens}

\subsection*{Your turn}
\begin{itemize}
\raggedright
  \item \emph{Say it:} \emph{próigrysh}
  \item \emph{Say it:} \emph{próigrysh}, once more
  \item \emph{Say it:} a sentence of your own with \emph{próigrysh}
  \item \emph{Recall:} say the Russian for a tie, then the Russian for a rival, then say \emph{próigrysh} again
\end{itemize}

\begin{tcolorbox}[breakable,colback=teal!4,colframe=teal!35!black,title={Before you move on}]
What does \ru{проигрыш} mean? (\textquotedblleft{}A loss (in a game)\textquotedblright{}.) Say it once more.
\end{tcolorbox}
\section[\texorpdfstring{\ru{лидировать} (lidírovat)}{лидировать (lidírovat)}]{\ru{лидировать} (lidírovat) — to be in the lead}

\label{lesson:RU-C510-lidirovat}



\begin{quote}
Before the new one: say the Russian for a rival, then the Russian for a loss.
\end{quote}

\subsection*{You'll want to know: \ru{лидировать}}
\textbf{\ru{лидировать}} (\emph{lidírovat}) — \textquotedblleft{}to be in the lead\textquotedblright{}.

\begin{quote}
\ru{После} \ru{первого} \ru{тайма} \ru{лидирует} \ru{наша} \ru{команда}.
\end{quote}

\begin{quote}
\emph{Pósle pérvogo táyma lidíruyet násha kománda.}
\end{quote}

\begin{quote}
After the first half our team is in the lead.
\end{quote}

Imperfective; present \ru{лидирую}. Built on \emph{\ru{лидер}} — a leader.

\begin{cousinweb}
\emph{\ru{лидировать}} is built on \emph{\ru{лидер}} 'a leader', which was borrowed from English \emph{leader} (EtymDB), with \emph{-\ru{ировать}}; EtymDB has no entry for the verb.
\end{cousinweb}

\begin{grammarlens}[title={What you've built}]
One more word for telling, explaining and describing at B1.
\end{grammarlens}

\subsection*{Your turn}
\begin{itemize}
\raggedright
  \item \emph{Say it:} \emph{lidírovat}
  \item \emph{Say it:} \emph{lidírovat}, once more
  \item \emph{Say it:} a sentence of your own with \emph{lidírovat}
  \item \emph{Recall:} say the Russian for a rival, then the Russian for a loss, then say \emph{lidírovat} again
\end{itemize}

\begin{tcolorbox}[breakable,colback=teal!4,colframe=teal!35!black,title={Before you move on}]
What does \ru{лидировать} mean? (\textquotedblleft{}To be in the lead\textquotedblright{}.) Say it once more.
\end{tcolorbox}
